\chapter{Error boundaries and recovery}
\label{ch:errorboundary2}

\begin{aside}
A widget that throws shouldn't take down the whole app. An
error boundary catches the throw and renders a fallback.
\end{aside}

\section{The boundary}

\begin{lstlisting}
auto safe_panel = error_boundary(
    [&] { return my_panel(); },
    [&](std::exception_ptr e) {
        return text("Error: " + to_string(e))
             | fg<theme::error>;
    }
);
\end{lstlisting}

The first lambda renders; if it throws, the second renders
instead.

\section{Catching}

\maya{}'s renderer wraps each error boundary's children in
a try/catch. Exceptions from layout, paint, or event
handlers are caught.

\section{Retry}

A fallback can include a ``retry'' button that re-runs the
failed render:

\begin{lstlisting}
return column(
    text("Failed: " + msg),
    button("Retry", [&] { reset_signal(...); })
);
\end{lstlisting}

\section{Reporting}

A boundary can report the error to the logger or to a
crash service.

\section{Scope}

Place boundaries at meaningful boundaries: per route, per
panel, per major feature. A boundary at the root catches
everything but loses context.

\section{What doesn't catch}

Exceptions in async tasks. The task system must wrap them
separately:

\begin{lstlisting}
ctx.async_effect([&]() -> Task<void> {
    try { co_await ...; }
    catch (...) { ctx.set(error, ...); }
});
\end{lstlisting}

\section{Recap}

\begin{recap}
\begin{itemize}[leftmargin=*]
  \item Error boundary catches render throws.
  \item Fallback renders alternative content.
  \item Include retry where reasonable.
  \item Async needs its own try/catch.
\end{itemize}
\end{recap}

\section*{Workshop}

\begin{enumerate}[leftmargin=*]
  \item Wrap each route in an error boundary.
  \item Add retry buttons.
  \item Wire boundary errors to your logger.
\end{enumerate}
